\documentclass[../../../ps_a.tex]{subfiles}

\begin{document}
    \begin{enumerate}
        \item%
            \hspace{-\zw}%
            \begin{minipage}[t]{\linewidth}
                \vspace{0pt}
                \centering
                    \includegraphics{\subfix{fig_em_ef_05/fig_em_ef_05_01_a.pdf}}%
            \end{minipage}
        \vspace{1pt}
        \item%
            $x$軸上での電位は$\phi = ax^2$である．%
            $x>0$の領域での電場は \:\,\uwave{$x$軸の負の向き} であり，強さは，
            
            \myspaceb%
            $\displaystyle%
                \left| E_x \right| = \left| \lim_{\Delta x \to 0}
                \frac{\Delta \phi}{\Delta x} \right| = 2ax \;.
            $\\
            $x=1\:\mathrm{m}$の位置では，

            \myspaceb%
            $\displaystyle%
                \left| E_x \right| = \uwave{20\:\mathrm{V/m}} \;.
            $
        \item%
            $y$軸上での電位は$\phi = -ay^2$である．%
            位置エネルギーが最大（運動エネルギーが最小）となる原点Oを通過できればよい．%
            原点に到達したときの速さを$u$とすれば，力学的エネルギー保存則より，
            
            \myspaceb%
            $\displaystyle%
                \myfraca{1}{2}mu^2 + q \cdot 0 = \myfraca{1}{2}m{v_0}^2 + q(-ah^2) \;.
            $\\
            よって，$\myfraca{1}{2}mu^2 = \myfraca{1}{2}m{v_0}^2 -qah^2 > 0$となるために，
            
            \myspaceb%
            $\displaystyle%
                \uwave{v_0 > h \sqrt{\frac{2qa}{m}}} \;.
            $
    \end{enumerate}
\end{document}

